\documentclass{article}

% packages
\usepackage{amsmath, amsthm, thmtools, amsfonts, amssymb, luacode, catchfile, tikzducks, hyperref, ifthen}
\ifcsname c@kobocompile\endcsname
	\usepackage[a5paper, total={1072pt, 1448pt}, margin=10pt, includeheadfoot]{geometry} % set page margins
\else
	\usepackage[a4paper, margin=50pt, includeheadfoot]{geometry}
\fi
\usepackage[shortlabels]{enumitem}
\usepackage[skip=3pt, indent=0pt]{parskip}

% language
\usepackage[bidi=basic, layout=tabular, provide=*]{babel}
\ifcsname c@english\endcsname
	\babelprovide[main, import]{english}
\else
	\babelprovide[main, import]{hebrew}
	\babelprovide{rl}
\fi
%\babelfont{rm}{Libertinus Serif}
\babelfont{rm}[Renderer=Harfbuzz]{Libertinus Serif}
\babelfont{sf}{Libertinus Sans}
\babelfont{tt}{Libertinus Mono}

% style
\AddToHook{cmd/section/before}{\clearpage}	% Add line break before section
\linespread{1.3}
\setcounter{secnumdepth}{0}		% Remove default number tags from sections, this won't do well with theorems
\AtBeginDocument{\setlength{\belowdisplayskip}{3pt}}
\AtBeginDocument{\setlength{\abovedisplayskip}{3pt}}
\graphicspath{ {../images/} }

% operators
\DeclareMathOperator\cis{cis}
\DeclareMathOperator\Sp{Sp}
\DeclareMathOperator\tr{tr}
\DeclareMathOperator\im{Im}
\DeclareMathOperator\re{Re}
\DeclareMathOperator\diag{diag}
\DeclareMathOperator*\lowlim{\underline{lim}}
\DeclareMathOperator*\uplim{\overline{lim}}
\DeclareMathOperator\rng{rng}
\DeclareMathOperator\Sym{Sym}
\DeclareMathOperator\Arg{Arg}
\DeclareMathOperator\Log{Log}
\DeclareMathOperator\dom{dom}
\DeclareMathOperator\supp{Supp}
\DeclareMathOperator\var{Var}
\DeclareMathOperator\cov{Cov}

% commands
%\renewcommand\qedsymbol{\textbf{מש''ל}}
%\renewcommand\qedsymbol{\fbox{\emoji{lizard}}}
\newcommand{\Aa}[0]{\mathcal{A}}
\newcommand{\Bb}[0]{\mathcal{B}}
\newcommand{\CC}[0]{\mathbb{C}}
\newcommand{\Cc}[0]{\mathcal{C}}
\newcommand{\EE}[0]{\mathbb{E}}
\newcommand{\FF}[0]{\mathbb{F}}
\newcommand{\Ff}[0]{\mathcal{F}}
\newcommand{\Ii}[0]{\mathcal{I}}
\newcommand{\Gg}[0]{\mathcal{G}}
\newcommand{\Ll}[0]{\mathcal{L}}
\newcommand{\Mm}[0]{\mathcal{M}}
\newcommand{\NN}[0]{\mathbb{N}}
\newcommand{\Nn}[0]{\mathcal{N}}
\newcommand{\PP}[0]{\mathbb{P}}
\newcommand{\Pp}[0]{\mathcal{P}}
\newcommand{\QQ}[0]{\mathbb{Q}}
\newcommand{\RR}[0]{\mathbb{R}}
\newcommand{\Rr}[0]{\mathcal{R}}
\newcommand{\Ss}[0]{\mathcal{S}}
\newcommand{\TT}[0]{\mathbb{T}}
\newcommand{\Uu}[0]{\mathcal{U}}
\newcommand{\Vv}[0]{\mathcal{V}}
\newcommand{\Ww}[0]{\mathcal{W}}
\newcommand{\ZZ}[0]{\mathbb{Z}}
\newcommand{\acts}[0]{\circlearrowright}
\newcommand{\explain}[2] {
	\begin{flalign*}
		 && \text{#2} && \text{#1}
	\end{flalign*}
}
\newcommand{\maketitleprint}[0]{ \begin{center}
	%\begin{tikzpicture}[scale=3]
	%	\duck[graduate=gray!20!black, tassel=red!70!black]
	%\end{tikzpicture}	
	\includegraphics[width=6cm]{cover}
\end{center}
}

% theorem commands
\newtheoremstyle{c_remark}
	{}	% Space above
	{}	% Space below
	{}% Body font
	{}	% Indent amount
	{\bfseries}	% Theorem head font
	{}	% Punctuation after theorem head
	{.5em}	% Space after theorem head
	{\thmname{#1}\thmnumber{ #2}\thmnote{ \normalfont{\text{(#3)}}}}	% head content
\newtheoremstyle{c_definition}
	{3pt}	% Space above
	{3pt}	% Space below
	{}% Body font
	{}	% Indent amount
	{\bfseries}	% Theorem head font
	{}	% Punctuation after theorem head
	{.5em}	% Space after theorem head
	{\thmname{#1}\thmnumber{ #2}\thmnote{ \normalfont{\text{(#3)}}}}	% head content
\newtheoremstyle{c_plain}
	{3pt}	% Space above
	{3pt}	% Space below
	{\itshape}% Body font
	{}	% Indent amount
	{\bfseries}	% Theorem head font
	{}	% Punctuation after theorem head
	{.5em}	% Space after theorem head
	{\thmname{#1}\thmnumber{ #2}\thmnote{ \text{(#3)}}}	% head content

\ifcsname c@english\endcsname
	\theoremstyle{plain}
	\newtheorem{theorem}{Theorem}[section]
	\newtheorem{lemma}[theorem]{Lemma}
	\newtheorem{proposition}[theorem]{Proposition}
	\newtheorem*{proposition*}{Proposition}
	%\newtheorem{corollary}[theorem]{אין חלופה עברית}

	\theoremstyle{definition}
	\newtheorem{definition}[theorem]{Definition}
	\newtheorem*{definition*}{Definition}
	\newtheorem{example}{Example}[section]
	\newtheorem{exercise}{Exercise}[section]

	\theoremstyle{remark}
	\newtheorem*{remark}{Remark}
	\newtheorem*{solution}{Solution}
	\newtheorem{conclusion}[theorem]{Conclusion}
	\newtheorem{notation}[theorem]{Notation}
\else
	\theoremstyle{c_plain}
	\newtheorem{theorem}{משפט}[section]
	\newtheorem{lemma}[theorem]{למה}
	\newtheorem{proposition}[theorem]{טענה}
	\newtheorem*{proposition*}{טענה}
	%\newtheorem{corollary}[theorem]{אין חלופה עברית}

	\theoremstyle{c_definition}
	\newtheorem{definition}[theorem]{הגדרה}
	\newtheorem*{definition*}{הגדרה}
	\newtheorem{example}{דוגמה}[section]
	\newtheorem{exercise}{תרגיל}[section]

	\theoremstyle{c_remark}
	\newtheorem*{remark}{הערה}
	\newtheorem*{solution}{פתרון}
	\newtheorem{conclusion}[theorem]{מסקנה}
	\newtheorem{notation}[theorem]{סימון}
\fi

% Questions related commands
\newcounter{question}
\setcounter{question}{1}
\newcounter{sub_question}
\setcounter{sub_question}{1}

\ifcsname c@english\endcsname
	\newcommand{\question}[1][0]{
		\ifthenelse{#1 = 0}{}{\setcounter{question}{#1}}
		\section{Question \arabic{question}}
		\addtocounter{question}{1}
		\setcounter{sub_question}{1}
	}

	\newcommand{\subquestion}[1][0]{
		\ifthenelse{#1 = 0}{}{\setcounter{sub_question}{#1}}
		\subsection{Part \alph{sub_question}}
		\addtocounter{sub_question}{1}
	}
\else
	\newcommand{\question}[1][0]{
		\ifthenelse{#1 = 0}{}{\setcounter{question}{#1}}
		\section{שאלה \arabic{question}}
		\addtocounter{question}{1}
		\setcounter{sub_question}{1}
	}

	\newcommand{\subquestion}[1][0]{
		\ifthenelse{#1 = 0}{}{\setcounter{sub_question}{#1}}
		\subsection{סעיף \localecounter{letters.gershayim}{sub_question}}
		\addtocounter{sub_question}{1}
	}
\fi

% import lua and start of document
\directlua{common = require ('../common')}

\GetEnv{AUTHOR}

% headers
\author{\AUTHOR}
\date\today

\title{פתרון ממ''ן 11 – חשבון אינפיניטסימלי 1 (20474)}

\begin{document}
\maketitle
\section{שאלה 1}
\subsection{סעיף א'}
נוכיח כי $a = k + m \sqrt{2}$ כאשר $k, m \in \NN$ הוא מספר אי־רציונלי. \\*
משפט 1.21 גורס כי $\sqrt{2}$ איננו רציונלי, לכן לא קיימים מספרים $p, q \in \ZZ$ כך ש־$\sqrt{2} = \frac{p}{q}$.
אילו היה המספר $m \sqrt{2}$ מספר רציונלי אז היה מתקיים $m \sqrt{2} = \frac{mp}{q}$ בניגוד לאי־רציונליות $\sqrt{2}$,
לכן $m\sqrt{2}$ מספר אי־רציונלי.
נגדיר $p', q' \in \ZZ$, מתקיים:
\[
	\frac{p}{q} + \frac{p'}{q'} = \frac{p q' + p'q}{q q'}
\]
לכן מספר כל חיבור מספרים רציונליים הוא מספר רציונלי, אילו הייתה התוצאה $k + m\sqrt{2}$ רציונלית אז היה נגזר כי גם $m\sqrt{2}$ רציונלית,
בסתירה להיותו אי־רציונלי. לכן המספר $a$ אי־רציונלי.

\subsection{סעיף ב'}
נוכיח כי לכל $n \in \NN$ המספר ${(1 + \sqrt{2})}^n$ לא רציונלי באינדוקציה: \\*
\textbf{בסיס האינדוקציה}:
המספר ${(1 + \sqrt{2})}^1$ הוא כמובן מספר אי־רציונלי בהתאם לסעיף א'. \\*
\textbf{מהלך האינדוקציה}:
נניח כי ${(1 + \sqrt{2})}^n = k + m\sqrt{2}, k, m \in \NN$ אי־רציונלי. מתקיים
\[
	{(1 + \sqrt{2})}^{n + 1}
	= {(1 + \sqrt{2})}^n {(1 + \sqrt{2})}
	= (k + m \sqrt{2}) (1 + \sqrt{2})
	= k + m \sqrt{2} + k \sqrt{2} + 2m
	= (k + 2m) + (m + k) \sqrt{2}
\]
ידוע כי $k + 2m, m + k \in \NN$ ולכן ${(1 + \sqrt{2})}^{n + 1}$ הוא מספר אי־רציונלי על־פי סעיף א' והושלם מהלך האינדוקציה.

\section{שאלה 2}
\subsection{סעיף א'}
יהיו $a, b \in \RR$ נוכיח שמתקיים:
\[
	\left| \sqrt{\left| a \right| + 1} - \sqrt{\left| b \right| + 1} \right| \le \frac{|a - b|}{2}
\]
מתקיים לכל $x, y \ge 0$ כאשר $x, y \in \RR$:
\[
	\sqrt{x} - \sqrt{y}
	= (\sqrt{x} - \sqrt{y}) \frac{(\sqrt{x} + \sqrt{y})}{(\sqrt{x} + \sqrt{y})}
	= \frac{(\sqrt{x} - \sqrt{y}) (\sqrt{x} + \sqrt{y})}{(\sqrt{x} + \sqrt{y})}
	= \frac{x - y}{\sqrt{x} + \sqrt{y}}
\]
נגדיר
\[
	x = |a| + 1, y = |b| + 1
\]
אז בשל חיוביות הערך המוחלט מתקיים לכל $a, b \in \RR$:
\[
	\left| \sqrt{\left| a \right| + 1} - \sqrt{\left| b \right| + 1} \right|
	= \frac{|a| + 1 - |b| - 1}{\sqrt{\left| a \right| + 1} + \sqrt{\left| b \right| + 1}}
	= \frac{|a| - |b|}{\sqrt{\left| a \right| + 1} + \sqrt{\left| b \right| + 1}}
\]
מהגדרת הערך המוחלט נובע כי לכל $a \in \RR$:
\[
	0 \le |a| \rightarrow 1 \le |a| + 1 \rightarrow \sqrt{1} = 1 \le \sqrt{|a| + 1}
\]
ולכן תמיד מתקיים
\[
	2 \le \sqrt{\left| a \right| + 1} + \sqrt{\left| b \right| + 1}
\]
בשל הופעת הביטוי במכנה נוכל להניח כי:
\[
	\left| \sqrt{\left| a \right| + 1} - \sqrt{\left| b \right| + 1} \right| = 
	\frac{|a| - |b|}{\sqrt{\left| a \right| + 1} + \sqrt{\left| b \right| + 1}}
	\ge \frac{|a| - |b|}{2}
\]
ועל־פי טענה 1.51:
\[
	\left| \sqrt{\left| a \right| + 1} - \sqrt{\left| b \right| + 1} \right|
	\ge \frac{|a - b|}{2}
\]

\subsection{סעיף ב'}
נוכיח
\[
	{\left( \frac{a + |a|}{2} \right)}^2
	+ {\left( \frac{a - |a|}{2} \right)}^2
	= a^2
\]
לכל מספר ממשי $n$ מתקיים $a^2 = |a| \cdot |a|$
לפי טענה 1.48 סעיף 6. נשתמש בטענה זו ונפתח סוגריים:
\[
	{\left( \frac{a + |a|}{2} \right)}^2
	+ {\left( \frac{a - |a|}{2} \right)}^2
	=
	\frac{a^2 + 2 a |a| + a^2}{4}
	+ \frac{a^2 - 2 a |a| + a^2}{4}
	=
	\frac{4 a^2 + 2 a |a| - 2 a |a|}{4}
	= a^2
\]

\section{שאלה 3}
\subsection{סעיף א'}
נוכיח כי לכל $x, y \in \RR$ אם $x \le y$ אז $\lfloor x \rfloor \le \lfloor y \rfloor$. \\*
נניח כי $x \le y$. נגדיר $x' = \langle x \rangle, y' = \langle y \rangle$ החלקים השבריים של $x, y$ בהאתמה.
מתקיים $x = \lfloor x \rfloor + x', y = \lfloor y \rfloor + y'$ לכן
\[
	\lfloor x \rfloor + x' \le \lfloor y \rfloor + y'
\]
אילו $x' \le y'$ הטענה מתקיימת, לכן נניח $x' > y'$.
מתקיים $\lfloor x \rfloor < \lfloor y \rfloor$, אחרת יתכן כי החלק השלם של המספרים יהיה זהה, אך החלק השברי $x'$ יהיה גדול מ־$y'$ ויגרור $y < x$ בסתירה להנחה.
אז מתקיים בכל מקרה $\lfloor x \rfloor \le \lfloor y \rfloor$.

\subsection{סעיף ב'}
\textbf{(i)}
נפתור את המשוואה
\begin{align*}
	& \left\lfloor x - \frac{1}{2} \right\rfloor^2 = 25 \\
	& \left(\left\lfloor x - \frac{1}{2} \right\rfloor^2 + 5\right)\left(\left\lfloor x - \frac{1}{2} \right\rfloor^2 - 5\right) = 0 \\
	& \left\lfloor x - \frac{1}{2} \right\rfloor = \pm 5
\end{align*}

על־פי טענה 1.64(1)
\begin{align*}
	& \pm 5 \le x - \frac{1}{2} < \pm 5 + 1 \\
	& \pm 5 + \frac{1}{2} \le x < \pm 5 + 1\frac{1}{2}  \\
	& -4 \frac{1}{2} \le x < -3\frac{1}{2}, 5\frac{1}{2} \le x < 6\frac{1}{2}
\end{align*}
\textbf{(ii)}
\begin{align*}
	& \lfloor x^2 \rfloor = 9 \\
	& 9 \le x^2 < 10 \\
	& 9 \le x^2 \land x^2 < 10 \\
	& 0 \le x^2 - 9 \land x^2 - 10 < 0 \\
	& 0 \le (x - 3)(x + 3) \land (x - \sqrt{10})(x + \sqrt{10}) < 0 \\
	& (x \ge 3 \lor x \le -3) \land (x < \sqrt{10} \lor x > -\sqrt{10}) \\
	& 3 \le x < \sqrt{10}, -\sqrt{10} < x \le -3
\end{align*}

\section{שאלה 4}
\subsection{סעיף א'}
תהי הקבוצה
\[
	A = \left\{ q \sqrt{3} \mid q \in (0, \infty) \cap \QQ \right\}
\]
נוכיח כי $A$ צפופה בקטע $[0, 1]$. \\*
ידוע כי $\QQ$ צפופה ב־$[0, 1]$.
נגדיר $f : \QQ \to A$ המוגדרת $f(x) = x \sqrt{3}$. ניתן לראות כי $f$ על וחד־חד ערכית ולכן הפיכה.
לכל שני מספרים $a, b \in A$ אשר מקיימים $a < b$ מתקיים גם $f^{-1}(x) < f^{-1}(y)$, והרי אלו הם מספרים רציונליים אשר ידוע כי צפופים,
נבחר מספר $c$ אשר מקיים $f^{-1}(x) < c < f^{-1}(y)$, ונשים לב כי גם מתקיים $x < f(c) < y$.
מצאנו כי $A$ צפופה לכל שני מספרים שונים המוכלים בה, לכן היא גם צפופה בתחום $[0, 1]$.

\subsection{סעיף ב'}
ננסח הגדרה לחוסר רציפות של קבוצה $A$ בקטע $I$.
תחילה נצרין את הטענה:
\[
	\lnot \forall x, y \in I
	(x < y \rightarrow \exists a \in A (x < a \land a < y)) 
\]
נפשט את הפסוק:
\[
	\exists x, y \in I
	(x < y \land \forall a \in A (x \ge a \lor a \ge y))
\]
ננסח את הפסוק שהתקבל: \\*
קבוצה $A$ תיקרא לא צפופה ב־$I$ אם קיימים שני מספרים $x, y \in I$
כך ש־$x < y$ וגם לכל איבר $a$ ב־$A$ מתקיים ש־$a \le x$ או ש־$a \ge y$.

\subsection{סעיף ג'}
נוכיח כי קבוצת השברים העשרוניים הסופיים שלא מופיעה בהם הספרה 3 אינה צפופה בקטע $[-1, 1]$. \\*
נגדיר $a = 0.3, b = 0.4$, ברור כי $a, b \in [-1, 1], a < b$, אך לכל מספר $x \in A$ מתקיים או $x \le a$ או $x \ge b$,
שכן כל מספר שלא עומד בהגדרה זו חייב להיות מהתצורה $0.3\ldots$ ולכן איננו שייך ל־$A$.
אז לפי הגדרת סעיף ב' הקבוצה $A$ איננה צפופה ב־$[-1, 1]$.
\end{document}
